\chapter{Gromov-Wasserstein Alignment}
\label{ch:gw_alignment}

While the supervised CCA baseline demonstrated that a linear cross-modal
mapping is attainable, relying on large-scale paired datasets fundamentally
limits the scalability and generalization of multimodal learning. This chapter
investigates whether the 3D and textual latent spaces can be aligned in a
purely unsupervised or semi-supervised manner by leveraging the mathematical
framework of Optimal Transport (OT).

Standard OT methods, such as the classic Wasserstein distance, require both
distributions to lie within the same metric space to compute a direct transport
cost. Because our 3D geometries and text captions inhabit completely disjoint
and heterogeneous dimensional spaces, we frame the alignment as a metric
measure problem using the Gromov-Wasserstein ($\mathcal{GW}$) distance.
Rather than comparing feature coordinates directly, $\mathcal{GW}$ attempts
to match the intra-modal topological structures---essentially aligning the
pairwise distance matrices of the two spaces. However, transitioning from
this elegant mathematical theory to empirical application requires navigating
severe computational bottlenecks and strict geometric assumptions. This chapter
details our implementation of scalable $\mathcal{GW}$ variants, documents a
comprehensive ablation study, and ultimately exposes the critical failure
points of the structural isometry assumption in real-world neural
representations.

% ─────────────────────────────────────────────────────────────────────────────
\section{Methods and Implementation}
\label{sec:gw_methods}

The standard $\mathcal{GW}$ optimization problem is notoriously expensive.
Solving for the exact dense coupling matrix $\pi$ generally scales at
$\mathcal{O}(n^3)$ with respect to the number of samples, while the memory
complexity for storing the cost matrices scales at $\mathcal{O}(n^2)$. Given
our objective to align tens of thousands of high-dimensional embeddings
extracted from the Objaverse dataset, computing the exact transport plan is
computationally intractable, even when utilizing the University of Padova's
HPC cluster infrastructure. To make this alignment feasible, we implemented
several advanced relaxations and hardware-accelerated OT solvers.

\subsection{Sliced GW and Rotation-Invariant Variants}
\label{subsec:sliced_gw}

One of the most persistent issues when processing diverse 3D objects from
expansive, crowdsourced datasets like Objaverse is the lack of canonical
orientation. Unlike carefully curated synthetic datasets, a vehicle in
Objaverse might be aligned along the positive Z-axis in one mesh and the
negative X-axis in another. If the geometric encoder does not inherently
guarantee rotation invariance, these arbitrary spatial orientations translate
directly into distorted intra-modal distance matrices, artificially inflating
the structural discrepancy between the 3D and text manifolds.

To mitigate this without forcing a computationally heavy pre-alignment phase,
we integrated rotation-invariant formulations of the Gromov-Wasserstein
problem. To further accelerate the optimization process, we employed Sliced
Gromov-Wasserstein (SGW) approximations \cite{vayer2019}. Instead of solving
the complex assignment problem in high-dimensional space, SGW iteratively
projects the feature distributions onto random 1-dimensional lines. In 1D
space, the optimal transport plan can be computed extremely efficiently in
$\mathcal{O}(n \log n)$ time using simple sorting algorithms. By integrating
these 1D costs over hundreds of random projection directions, we obtained a
robust approximation of the true $\mathcal{GW}$ distance, which allowed us
to aggressively scale our batch sizes during the initial exploratory
experiments without exhausting the memory limits of our NVIDIA A100 GPUs.

The Rotation-Invariant SGW (RISGW) variant extends this by explicitly
optimising a rotation matrix $\Delta$ on the Stiefel manifold (via Pymanopt
and JAX autodiff) to find the projection directions that minimise the SGW
cost. A Semi-Supervised extension (SS-RISGW) further incorporates an anchor
supervision penalty term, weighted by $\lambda \in \{0, 0.1, 1.0, 10.0\}$,
to weakly guide the rotation towards known correspondences.

\subsection{Fused GW and Anchor-Constrained FGW}
\label{subsec:fused_gw}

While pure $\mathcal{GW}$ relies exclusively on the topological structure
(the relational distance matrices $C_X$ and $C_Y$), Fused Gromov-Wasserstein
(FGW) \cite{vayer2019fgw} introduces a mechanism to combine structural
matching with direct feature matching. The FGW formulation utilizes a
hyperparameter $\alpha \in [0, 1]$ that dictates the trade-off between the
Wasserstein cost (feature alignment) and the Gromov-Wasserstein cost
(structural alignment):
\begin{equation}
    \text{FGW}(X, Y) = (1 - \alpha)\,W(X, Y) + \alpha\,\mathcal{GW}(X, Y)
    \label{eq:fgw}
\end{equation}

Because our modalities occupy disjoint spaces, a direct feature cost is only
possible after projecting both spaces into a shared CCA subspace. We therefore
operated FGW in the $k=50$ CCA subspace and swept $\alpha$ from 0 to 1,
finding the best result at $\alpha = 0.2$ with 1,000 anchor pairs.

We also evaluated an Anchor-Constrained FGW variant, in which the cost matrix
was modified to assign zero cost to the 1,000 known anchor pairs, with an
additional penalty on off-diagonal anchor entries (penalty values
$\in \{0, 10, 100\}$). The intent was to force the solver to respect known
ground-truth correspondences as fixed structural scaffolding. In both cases,
results were identical to standard FGW at $4.4\%$ Top-5, indicating that the
structural incompatibility of the spaces is too severe for even locally
supervised anchoring to overcome.

\subsection{Low-Rank GW via OTT-JAX}
\label{subsec:low_rank_gw}

For our definitive, large-scale unsupervised experiments, we transitioned our
pipeline to the OTT-JAX library. This framework is specifically built on top
of JAX to support hardware-accelerated, auto-differentiable optimal transport,
allowing us to fully exploit the parallelisation capabilities of the Slurm
cluster.

Rather than computing the full, dense transport plan
$\pi \in \mathbb{R}^{n \times n}$, we utilised the library's Low-Rank
Gromov-Wasserstein solver. This method factors the coupling matrix into two
much smaller matrices, effectively bounding the rank to a user-defined
parameter $r$ (where $r \ll n$). Factoring the transport plan not only
reduced the spatial complexity from quadratic to linear---enabling the
simultaneous processing of massively larger subsets of our 46,675 paired
samples---but it also served as an implicit regularizer. By restricting the
rank, the solver is forced to focus on aligning the dominant, macro-level
topological clusters of the latent spaces, intrinsically filtering out
high-frequency noise and minor geometric inconsistencies inherent to the
raw 3D mesh encodings.

% ─────────────────────────────────────────────────────────────────────────────
\section{Results}
\label{sec:gw_results}

Evaluating the quality of the alignments generated by the Gromov-Wasserstein
solvers requires a direct comparison against the supervised upper bound
established in Chapter~\ref{ch:baseline}. To ensure absolute fairness, all
OT-based matching and retrieval metrics were computed on the exact same
held-out query set of 500 pairs, utilising the distance representations
extracted from the PointNet and CLIP (ViT-bigG-14) backbones.

A preliminary glance at the quantitative outcomes reveals a stark contrast
between the theoretical elegance of optimal transport and its practical
applicability to complex, real-world neural representations. We dissect these
findings by first examining the purely unsupervised solvers, followed by an
analysis of the semi-supervised anchoring experiments.

\subsection{Quantitative Evaluation of Unsupervised GW}
\label{subsec:unsupervised_gw_results}

Our primary experimental objective was to achieve true zero-shot alignment
using Sliced Gromov-Wasserstein (SGW) and the Low-Rank Gromov-Wasserstein
(LR-GW) implementations. These solvers were tasked with constructing a
coupling matrix $\pi$ relying exclusively on the intra-modal pairwise
distance matrices $C_X$ and $C_Y$, without any cross-modal anchor points to
guide the optimisation. The retrieval accuracies are reported in
Table~\ref{tab:gw_results}.

\begin{table}[H]
\centering
\resizebox{\textwidth}{!}{%
\begin{tabular}{lcccccc}
\toprule
\textbf{Alignment Method} & \textbf{Sup.} & \textbf{Anc.}
    & \textbf{Match} & \textbf{Top-1} & \textbf{Top-5} & \textbf{Top-10} \\
\midrule
\multicolumn{7}{l}{\textit{Supervised Baseline (Upper Bound)}} \\
CCA + Affine ($k = 50$)            & \checkmark & 30k & 13.5\% & 14.6\% & \textbf{35.7\%} & 46.4\% \\
\midrule
\multicolumn{7}{l}{\textit{Fully Unsupervised Optimal Transport}} \\
Sliced GW (100 proj., $k$=20)      & $\times$   & 5k  & 1.5\%  & 1.1\%  & 3.2\%           & 4.8\%  \\
RISGW (Stiefel, 30 iter.)          & $\times$   & 1k  & 0.4\%  & 0.2\%  & 1.6\%           & 2.2\%  \\
Low-Rank GW ($r$=20, angular)      & $\times$   & 30k & 1.6\%  & ---    & 1.2\%           & ---    \\
\midrule
\multicolumn{7}{l}{\textit{Semi-Supervised / Constrained Optimal Transport}} \\
SS-RISGW ($\lambda$=1.0, best)     & $\sim$     & 1k  & 0.2\%  & 0.2\%  & 1.4\%           & 2.6\%  \\
FGW ($\alpha$=0.2, best)           & $\times$   & 1k  & 3.6\%  & 3.6\%  & \textbf{4.4\%}  & 5.4\%  \\
FGW anchor-constrained (pen.=10)   & $\sim$     & 1k  & 3.6\%  & 3.6\%  & 4.4\%           & 5.4\%  \\
\midrule
\multicolumn{7}{l}{\textit{Random Baseline}} \\
Independent Gaussian               & ---        & --- & $\approx$0.6\% & $\approx$0.6\%
                                                      & $\approx$1.6\% & $\approx$2.0\% \\
\bottomrule
\end{tabular}%
}
\caption{Retrieval performance of all Gromov-Wasserstein variants on
PointNet $\times$ CLIP. $\times$ = unsupervised, $\sim$ = partially
supervised, \checkmark = fully supervised. The best OT method (FGW,
4.4\% Top-5) is $8\times$ below CCA+Affine (35.7\%).
CCA+Affine result taken from Table~\ref{tab:baseline_results}
(Chapter~\ref{ch:baseline}); all other results: same protocol,
30,000 anchor pairs, 500 query points, $k=50$, averaged over 3 seeds.}
\label{tab:gw_results}
\end{table}

The performance of all purely unsupervised methods experienced a catastrophic
drop compared to the linear baseline. In several runs, the exact matching
accuracy hovered marginally above random chance. Rather than discarding this
as a mere algorithmic failure, inspecting the output of the optimal transport
plan provides deep insights into the latent space topologies.

Unlike the CCA projection---which forcibly aligns the axes of maximum variance
based on known semantic pairs---the GW solver operates purely on geometric
faith. It attempts to overlay the shape of the 3D distance manifold directly
onto the text distance manifold. The extreme degradation in Top-$K$ retrieval
indicates that minimizing the global metric measure discrepancy (the GW
objective) does not correlate with minimizing the semantic discrepancy. In
practice, the optimiser frequently fell into localised topological traps. For
instance, tightly clustered groups of 3D objects (such as the heavily
populated categories of ``furniture'' or ``vehicles'' in Objaverse) were
mapped en masse to completely unrelated but similarly dense clusters in the
text space. The algorithm successfully matched the statistical \textit{density}
and \textit{variance} of the distributions, but completely scrambled their
\textit{identities}.

Furthermore, the low-rank constraint in OTT-JAX, while computationally
necessary to avoid out-of-memory errors on the HPC cluster, inherently forced
the coupling matrix to smooth out high-frequency matching details. While this
filtering effect is beneficial for stabilising gradients during optimisation,
it proved detrimental for exact instance-level retrieval, often resulting in
``many-to-one'' assignments where multiple distinct 3D queries collapsed onto
a single, centrally located text embedding.

\subsection{The Impact of Semi-Supervised Anchoring (FGW)}
\label{subsec:fgw_results}

Given the collapse of the purely unsupervised approach, we analysed whether
injecting a partial supervised signal into the transport plan could rescue
the alignment. The Fused Gromov-Wasserstein (FGW) framework allows this by
combining a direct feature-distance term with the structural GW term via the
trade-off parameter $\alpha$. We tested $\alpha = 0.2$ on 1,000 anchor
pairs---a configuration that places stronger weight on the Wasserstein
(feature-distance) term relative to the structural GW term.

We also evaluated an Anchor-Constrained FGW variant, setting cost zero on
the 1,000 known anchor pairs in the feature-distance matrix, with an
additional penalty on off-diagonal anchor entries (penalty $\in \{0, 10,
100\}$). The intent was to use the known correspondences as immovable
structural pegs around which the remaining unsupervised assignments would
be constructed.

Both configurations produced identical results: 4.4\% Top-5, the same as
standard FGW (Table~\ref{tab:gw_results}). The anchor constraints failed to
propagate any additional alignment signal across the broader geometry. This
outcome strongly suggests that the structural incompatibility between the two
spaces is too fundamental for even locally supervised anchoring to overcome
--- a hypothesis we investigate rigorously in the ablation study of
Section~\ref{sec:ablation_study}.

% ─────────────────────────────────────────────────────────────────────────────
\section{Ablation Study --- Thirteen Diagnostic Tests}
\label{sec:ablation_study}

The catastrophic degradation of retrieval accuracy observed in our
unsupervised experiments necessitated a rigorous diagnostic phase. When an
optimisation algorithm fails to converge to a meaningful semantic mapping,
the failure can typically be attributed to one of three sources: a bug in
the code implementation, an unstable hyperparameter configuration, or a
fundamental incompatibility in the underlying mathematical assumptions of
the data.

To systematically isolate the root cause, we designed and executed an
exhaustive ablation study comprising thirteen distinct experimental
configurations. This rigorous testing framework was designed to challenge
every variable in our pipeline sequentially, ensuring that the poor
performance of Gromov-Wasserstein was a genuine reflection of the latent
space topologies rather than a technical oversight. Tests 1 through 11 cover
standard algorithmic variations, while Tests 12 and 13 address advanced
methodological verifications regarding feature preprocessing and entropic
smoothing, and were added following specific instructions from the supervisor.

\subsection{Implementation Validation (Tests 1 and 2)}
\label{subsec:implementation_validation}

Before questioning the complex geometry of the Objaverse dataset, it was
imperative to guarantee that our custom Optimal Transport solvers,
particularly the Low-Rank Gromov-Wasserstein (LR-GW) implementation via
OTT-JAX, were structurally sound.

\paragraph{Test 1: Synthetic Sanity Check.}
We constructed synthetic data scenarios where the optimal alignment was
mathematically guaranteed to exist. We defined a target space $Y = XQ$,
where $Q$ is a random orthogonal matrix representing a perfect geometric
isometry (a rotation in high-dimensional space). Operating with $n=500$,
$d=50$, and a rank constraint of 200, the LR-GW solver perfectly recovered
the mapping, achieving a $93.4\%$ Top-5 accuracy on the rotated data and
$91.6\%$ on reflected data ($Y = X$ with a sign flip). Conversely, when
attempting to map the space onto pure independent Gaussian noise, performance
immediately collapsed to a random-guessing baseline of $1.6\%$. This test
definitively proved that the solver is capable of finding isometric alignments
when they genuinely exist in the underlying data.

\begin{table}[htpb]
\centering
\small
\begin{tabular}{llcc}
\toprule
\textbf{Case} & \textbf{Construction} & \textbf{Match (\%)} & \textbf{Top-5 (\%)} \\
\midrule
A -- Rotation       & $Y = XQ$, perfect isometry      & 67.0 & 93.4 \\
B -- Reflection     & $Y = X$ with sign flip          & 75.2 & 91.6 \\
C -- Rotation+Noise & $Y = XQ + \mathcal{N}(0, 0.1)$ &  1.0 &  1.8 \\
D -- Pure Noise     & $Y$ independent Gaussian        &  0.6 &  1.6 \\
\bottomrule
\end{tabular}
\caption{Sanity check (Test 1), LR-GW ($r=200$) on synthetic data.
Cases A--B confirm the solver recovers true isometric mappings.
Case D is the random-guessing baseline.}
\label{tab:sanity_check}
\end{table}

\paragraph{Test 2: Self-GW on Real Features.}
To ensure the Objaverse features themselves did not contain numerical
anomalies preventing OT convergence, we tasked the solver with the trivial
objective of aligning a modality with itself. We computed the $\mathcal{GW}$
distance between the PointNet features and themselves, and identically for
the CLIP features. The algorithm achieved $86.6\%$ Top-5 retrieval for
self-aligned PointNet and $88.0\%$ Top-5 for self-aligned CLIP. Notably,
the structural cost of aligning the two distinct spaces (PointNet to CLIP)
was $1.4\times$ higher than the cost of aligning identical spaces. These
initial tests confirm the implementation is correct; the failure resides
within the cross-modal data topology.

\begin{table}[htpb]
\centering
\small
\begin{tabular}{lccc}
\toprule
\textbf{Comparison} & \textbf{GW Cost} & \textbf{Match (\%)} & \textbf{Top-5 (\%)} \\
\midrule
GW(PointNet, PointNet) -- self & 0.0233 & 91.4 & 86.6 \\
GW(CLIP, CLIP) -- self         & 0.0242 & 71.2 & 88.0 \\
GW(PointNet, CLIP) -- real     & 0.0332 &  0.2 &  0.6 \\
\bottomrule
\end{tabular}
\caption{Self-GW validation (Test 2), rank=200. The structural cost of
aligning distinct modalities (0.0332) is $1.4\times$ higher than
self-alignment, quantifying the geometric incompatibility of the spaces.}
\label{tab:self_gw}
\end{table}

\subsection{Hyperparameter and Metric Sweeps (Tests 3 to 6)}
\label{subsec:hyperparameter_sweeps}

Having validated the codebase, we turned our attention to the optimisation
dynamics. The $\mathcal{GW}$ loss landscape is highly non-convex, and
gradient-based solvers can easily become trapped in localised topological
artefacts if not properly regularised.

\paragraph{Test 3: Epsilon Regularisation Sweep.}
Entropic regularisation ($\epsilon$) dictates the smoothness of the transport
plan. We conducted a sweep across eight distinct values
($\epsilon \in \{0.001, 0.005, 0.01, 0.05, 0.1, 0.5, 1.0, 5.0\}$) using
rank 20 and cosine cost. The solver returned an identical final GW cost of
$0.0392$ across all eight configurations. This indicates that the algorithm
converges to the exact same stationary point regardless of the regularisation
strength, suggesting the absence of better local minima in the loss landscape.

\paragraph{Test 4: Cost Function Formulation.}
The geometry of high-dimensional neural representations often behaves
counter-intuitively under standard Euclidean norms. We replaced the default
squared-Euclidean distance with an angular distance metric
$d(x, x') = \arccos\!\left(\tfrac{x \cdot x'}{\|x\|\,\|x'\|}\right)$,
following the formulations of Li et al.~\cite{li2026}. While the angular
metric improved the Top-5 accuracy slightly from $0.4\%$ to $1.9\%$, the
result remained roughly $18\times$ below the supervised CCA+Affine baseline
\cite{hadgi2025}. The choice of distance function is not the limiting factor.

\paragraph{Test 5: Raw Features and Scale-Matching.}
We investigated whether normalisation techniques were at fault. Applying
LR-GW directly on raw, pre-CCA normalised features, we incorporated the
median-ratio scale correction $s = \text{med}(D_Y)/\text{med}(D_X)$ of
Li et al.~\cite{li2026} to adjust for modal density differences. This
attempt yielded a dismal $0.5\%$ Top-5 accuracy, confirming that simple
scale discrepancies are not the core issue.

\paragraph{Test 6: Rank Constraint Sweep.}
We aggressively swept the rank parameter across seven values:
$r \in \{2, 5, 10, 20, 50, 100, 200\}$. The optimal performance peaked at
rank 50, generating a Top-5 accuracy of $2.4\%$. In the synthetic Test 1,
raising the rank to 200 drastically improved accuracy (from $19\%$ to
$93\%$); on real data, no such improvement occurred. This proves the
bottleneck is not the model's representational capacity, but the structural
incompatibility of the data itself.

\subsection{Dataset and Architecture Generalisation (Tests 7 to 9)}
\label{subsec:dataset_generalization}

We further examined whether the constraints of our dataset sampling or the
specific architectures chosen were to blame for the alignment collapse.

\paragraph{Test 7: Query Set Expansion.}
It was hypothesised that a query set of 500 pairs might not provide enough
contextual density for the solver to orient the manifolds correctly. We
scaled the query size to $n \in \{500, 1000, 2000\}$. The alignment
performance did not improve with more data; in fact, it slightly degraded at
the maximum size of 2000 queries, proving that data scarcity is not the
limiting factor.

\paragraph{Test 8: CCA Dimension Sweep.}
We tested whether the CCA subspace size $k=50$ was discarding too much
structural variance. We evaluated $k \in \{10, 20, 50, 100\}$. The best
LR-GW result occurred at $k=20$, achieving $3.0\%$ Top-5. By comparison,
the supervised CCA+Affine baseline evaluated in that exact same $k=20$
subspace reached $30.1\%$. The massive performance gap persists consistently
across all dimensionalities.

\paragraph{Test 9: Comprehensive Encoder Matrix.}
To rule out a specific flaw in the PointNet or CLIP topologies, we repeated
the standard LR-GW configuration across all six encoder combinations
(PointNet and SparseConv against CLIP, RoBERTa, and BERT) across three
independent seeds. The unsupervised retrieval metrics remained uniformly
flat, hovering tightly between $1.1\%$ and $1.7\%$ Top-5 regardless of
the encoder combination, conclusively demonstrating that the failure is an
intrinsic property of cross-modal alignment rather than an encoder-specific
anomaly.

\subsection{Library Consistency and Transport Plan Analysis (Tests 10 and 11)}
\label{subsec:library_analysis}

\paragraph{Test 10: Cross-Library Verification.}
To mathematically insulate our findings from potential bugs hidden deep within
a specific software ecosystem, we replicated the core experiments using the
Python Optimal Transport (POT) library \cite{flamary2021pot}, which relies on
fundamentally different numerical solvers than OTT-JAX. The Sliced
Gromov-Wasserstein (SGW) implementation via POT generated a $3.4\%$ Top-5
accuracy, confirming that two independent, state-of-the-art optimal transport
libraries agree on the same magnitude of failure.

\paragraph{Test 11: Transport Plan Inspection.}
We directly analysed the raw structure of the generated transport matrix $T$.
A successful retrieval matrix heavily concentrates probabilistic mass along
its diagonal. Our analysis revealed that the LR-GW transport plan's
diagonal-to-uniform ratio languished between $0.4\times$ and $1.2\times$.
Furthermore, its mathematical entropy was measured between $7.4$ and $9.6$
(approaching the theoretical maximum of $\log(n^2) = 12.4$ for a fully
uniform plan). This statistically proves that the matrix is
indistinguishable from a uniform random distribution; the solver finds
absolutely no semantic signal in the cross-modal structure.

\subsection{Advanced Methodological Verifications (Tests 12 and 13)}
\label{subsec:advanced_verifications}

Following a preliminary review of Tests 1 through 11, two final experiments
were designed following specific instructions from the supervisor to
definitively close any remaining methodological loopholes.

\paragraph{Test 12: GW on Raw Features (Pre-CCA).}
A valid methodological concern was that the CCA projection, trained
explicitly on 30,000 anchor pairs, might linearly filter out complex
structural geometries that the non-linear $\mathcal{GW}$ could have
exploited. To test this, we applied the solver directly to the raw,
unprojected PointNet (1024-dim) and CLIP (1280-dim) features, entirely
bypassing CCA. LR-GW handles different input dimensions because it operates
on $n \times n$ internal distance matrices, not on the feature vectors
directly. We tested four normalisation variants to exclude any pre-processing
artefact:

\begin{table}[htpb]
\centering
\small
\begin{tabular}{lcc}
\toprule
\textbf{Pre-processing Variant} & \textbf{Match (\%)} & \textbf{Top-5 (\%)} \\
\midrule
$L_2$-norm, no scale-matching                        & 0.4 & 1.5 \\
$L_2$-norm + Li et al.\ scale-matching ($s$=28.25)  & 0.4 & 1.5 \\
StandardScaler (zero-mean, unit-variance), no SM     & 0.5 & 1.3 \\
StandardScaler + Li et al.\ scale-matching           & 0.6 & 1.5 \\
\bottomrule
\end{tabular}
\caption{LR-GW on raw pre-CCA features (Test 12). All normalisation
variants remain at or below random guessing ($\approx$1\%), proving the
alignment failure is not an artefact of the CCA projection.}
\label{tab:raw_features_test}
\end{table}

All four variants yielded Top-5 accuracies between $1.3\%$ and $1.5\%$,
statistically indistinguishable from random guessing. The scale factor
$s = 28.25$ confirms that the two raw feature spaces operate on very
different distance scales, yet even after correcting for this via
scale-matching, performance does not improve. This conclusively proves
that the supervised CCA step did not destroy exploitable structural
information --- the structural isometry simply never existed in the raw
features.

\paragraph{Test 13: Fine Epsilon Sweep in the Low Range.}
The final verification targeted the entropic smoothing factor. It was
hypothesised that standard $\epsilon$ values might be excessively smoothing
the landscape, masking a sharp, correct global minimum. We conducted a
fine-grained sweep spanning three orders of magnitude in the
near-unregularised regime:
$\epsilon \in \{10^{-4}, 5 \cdot 10^{-4}, 10^{-3}, 5 \cdot 10^{-3},
10^{-2}, 5 \cdot 10^{-2}, 10^{-1}\}$. The computed $\mathcal{GW}$ cost
remained perfectly static at $0.0392$ (post-CCA) and $0.4854$ (pre-CCA)
across all seven values, with zero variation up to four decimal places
($\Delta\text{cost} = 0.0000$). Reducing the smoothing penalty from $0.1$
to $0.0001$ had no effect on retrieval performance. This is the definitive
computational signature of a mathematically flat loss landscape: the solver
is not stuck due to excessive entropy, it converges to the same structural
mapping because no better semantic alignment is geometrically possible.

\begin{table}[H]
\centering
\resizebox{\textwidth}{!}{%
\begin{tabular}{clp{9cm}}
\toprule
\textbf{Test} & \textbf{Targeted Variable}
    & \textbf{Conclusion / Alternative Ruled Out} \\
\midrule
1  & Synthetic sanity check
    & Implementation bug (93.4\% Top-5 on true isometry \checkmark) \\
2  & Self-GW on real features
    & Data numerical issues (91.4\% match on self \checkmark) \\
3  & $\epsilon$ sweep ($0.001$--$5.0$, 8 values)
    & Wrong regularisation (cost 0.0392 identical for all values) \\
4  & Cost function (Euclidean vs.\ angular)
    & Wrong distance metric (angular: 1.9\%, $18\times$ below baseline) \\
5  & Raw features + scale-matching
    & Post-CCA distortion (raw features equally bad at 0.5\%) \\
6  & Rank sweep ($r \in \{2 \ldots 200\}$, 7 values)
    & Insufficient capacity (best rank=50 gives 2.4\%) \\
7  & Query set size ($500$--$2000$)
    & Data scarcity (more data does not help) \\
8  & CCA dimension ($k \in \{10 \ldots 100\}$)
    & Wrong subspace (best $k$=20 gives 3.0\% vs.\ 30.1\% for CCA+Affine) \\
9  & All 6 encoder pairs (3 seeds each)
    & Encoder-specific artefact (1.1--1.7\% Top-5 uniformly) \\
10 & POT vs.\ OTT-JAX libraries
    & Library bug (both independently agree on same failure) \\
11 & Transport plan structure
    & Hidden signal (plan entropy statistically identical to uniform random) \\
\midrule
\multicolumn{3}{c}{\textit{Advanced verifications --- added following
supervisor instructions}} \\
\midrule
12 & LR-GW on raw pre-CCA features (4 variants)
    & CCA filtering exploitable structure
      (pre-CCA identical to post-CCA at $\approx$1.5\%) \\
13 & Fine $\epsilon$ sweep ($10^{-4}$--$10^{-1}$, 7 values)
    & Excessive smoothing masking solution
      (cost perfectly stable to 4 decimal places) \\
\bottomrule
\end{tabular}%
}
\caption{All thirteen diagnostic ablation tests and the specific alternative
explanation each one definitively rules out.}
\label{tab:thirteen_tests_summary}
\end{table}

% ─────────────────────────────────────────────────────────────────────────────
\section{Theoretical Analysis: Why GW Cannot Work}
\label{sec:why_gw_fails}

The exhaustive ablation study presented in Section~\ref{sec:ablation_study}
systematically eliminated every conventional point of computational failure.
With the implementation, hyperparameters, and libraries fully validated, the
persistent inability of Gromov-Wasserstein to align 3D objects with textual
descriptions must be attributed to a fundamental violation of the mathematical
assumptions underlying the Optimal Transport framework.

This section formalises the theoretical reasons behind this collapse,
contrasting the intrinsic topological mismatch of the latent spaces against
the success of the supervised linear baseline.

\subsection{The Absence of Structural Isometry}
\label{subsec:isometry_bottleneck}

The theoretical justification for employing the Gromov-Wasserstein distance
is its ability to bypass heterogeneous coordinate systems by matching
relational intra-modal distance matrices. However, a foundational theorem by
Vayer et al.~\cite{vayer2019} proves that the GW distance between two spaces
converges to a meaningful zero-cost minimum \textit{if and only if} those
spaces are globally isometric --- meaning there exists a rigorous,
distance-preserving bijection between their structures.

Our empirical results definitively prove that the PointNet and CLIP latent
spaces do not share an isometric topology. This is theoretically corroborated
by Hadgi et al.~\cite{hadgi2025}, who evaluated the structural similarity of
these unaligned spaces using Linear Centered Kernel Alignment (CKA). While
standard functional vision-text models typically exhibit CKA scores between
$0.30$ and $0.48$, the unaligned PointNet and CLIP spaces yielded a CKA
score of just $0.12$, as reported in their Figure 2. This extremely low
correlation quantifies a profound structural disconnect: uni-modal 3D
encoders do not naturally converge to topological manifolds that mirror
textual representations.

Furthermore, applying the methodology of Li et al.~\cite{li2026} --- who
utilise GW strictly as a training-free thermometer for latent space
compatibility rather than an alignment tool --- our encoders consistently
converged to a GW cost of approximately $0.031$ to $0.032$ post-CCA.
According to their validated metric, which correlates strongly with
Vision-Language Model (VLM) downstream performance ($R^2 = 0.43$ across
$60+$ full training runs), this cost threshold explicitly flags our feature
spaces as an intrinsically low-compatibility setting. Without an underlying
structural homology, the GW solver has no geometric signal to exploit,
reducing the transport plan to a mathematically valid but semantically
meaningless mapping.

\subsection{Final Method Ranking and The Necessity of Supervision}
\label{subsec:necessity_supervision}

The fundamental dichotomy between the successful linear baseline and the
failed optimal transport solvers lies entirely in the exploitation of
ground-truth correspondence.

\begin{table}[H]
\centering
\resizebox{\textwidth}{!}{%
\begin{tabular}{llcccc}
\toprule
\textbf{Category} & \textbf{Method}
    & \textbf{Match} & \textbf{Top-1} & \textbf{Top-5} & \textbf{Top-10} \\
\midrule
\multicolumn{6}{l}{\textit{Unsupervised Optimal Transport}} \\
& Sliced GW ($k$=20, 5k anc.)       & 1.5\%  & 1.1\%  & 3.2\%           & 4.8\%  \\
& FGW ($\alpha$=0.2, 1k anc.)       & 3.6\%  & 3.6\%  & 4.4\%           & 5.4\%  \\
\midrule
\multicolumn{6}{l}{\textit{Supervised Linear}} \\
& CCA + Affine (30k anc.)           & 13.5\% & 14.6\% & \textbf{35.7\%} & 46.4\% \\
\midrule
\multicolumn{6}{l}{\textit{Supervised Non-Linear
    (Chapter~\ref{ch:nonlinear})}} \\
& KCCA Nystr\"{o}m ($\gamma$=0.1)   & 4.7\%  & 4.9\%  & 15.8\%          & 23.3\% \\
& Riemannian Metric Learning        & 9.4\%  & 8.7\%  & 27.2\%          & 38.4\% \\
\midrule
\multicolumn{6}{l}{\textit{Random Baseline}} \\
& Independent Gaussian              & $\approx$0.6\% & $\approx$0.6\%
                                    & $\approx$1.6\% & $\approx$2.0\% \\
\bottomrule
\end{tabular}%
}
\caption{Complete ranking of all tested alignment methods on PointNet
$\times$ CLIP. Final Top-5 order: CCA+Affine (35.7\%) $>$ RML (27.2\%)
$>$ KCCA (15.8\%) $>$ FGW (4.4\%) $>$ SGW (3.2\%) $\approx$ random
($\approx$1.6\%). Supervised methods consistently and substantially
outperform all unsupervised alternatives.
CCA+Affine result taken from Table~\ref{tab:baseline_results}
(Chapter~\ref{ch:baseline}). KCCA and RML results from
Chapter~\ref{ch:nonlinear}. All results: 30,000 anchor pairs,
500 query points, $k=50$, averaged over 3 random seeds.}
\label{tab:final_ranking}
\end{table}

Why does the CCA + Affine pipeline succeed where GW fails? The affine mapping
$T(x) = Rx + b$ strictly minimises the mathematical discrepancy between
predefined semantic anchor pairs ($\|Rx + b - y\|^2$ directly over known
pairs), completely overriding the natural global topology of the manifolds.
It does not require the 3D space to inherently resemble the text space
geometrically; it artificially forces the correlation through supervised
constraints.

Conversely, Gromov-Wasserstein operates blindly without labels. It can only
exploit structural similarities, which, as demonstrated by the CKA scores and
our thirteen ablation tests, are fundamentally absent between raw 3D
coordinate representations and linguistic transformer embeddings.

Attempting to bridge the semantic gap between non-isometric modalities using
purely unsupervised structural alignment is a theoretical dead end. Explicit
semantic supervision remains an absolute mathematical necessity for functional
3D-to-text retrieval in these specific latent spaces. This conclusion is
further reinforced by the non-linear supervised methods presented in
Chapter~\ref{ch:nonlinear}: even the most expressive non-linear alternatives
(RML at 27.2\%, KCCA at 15.8\%) fall short of the simple linear affine map
(35.7\%), confirming that the performance ceiling in this problem is set by
the quality of the supervision, not by the complexity of the alignment model.